\input{../tex-inputs/latex-leseansicht-vorspann}

\section[Arthur Schnitzler an Gustav Schwarzkopf, {[}zwischen 11. und 18. 12. 1897?{]}]{L04186 Arthur Schnitzler an Gustav Schwarzkopf, {[}zwischen 11. und 18. 12. 1897?{]}}
\nopagebreak\mylabel{L04186v}
\rehead{ }\normalsize\beginnumbering\briefempfaengerindex{Schwarzkopf, Gustav@\textsc{Schwarzkopf, Gustav}!zzzSchnitzler, Arthur@\emph{von Arthur Schnitzler}!1897-12-112@{{[}zwischen 11. und 18. 12. 1897?{]}}|(be}
\toendnotes[C]{\smallbreak\pagebreak[2]}
\correspDesc{Versand  durch Arthur Schnitzler im Zeitraum [zwischen 11. und 18. 12. 1897?] in Frankgasse 1
\newline{}Erhalt  durch Gustav Schwarzkopf im Zeitraum [zwischen 11. und 18. 12. 1897?] in Tiefer Graben 23}\toendnotes[C]{\smallbreak}
\Standort{CUL, Schnitzler, B 96.}
\physDesc{,  Blätter,  Seiten, 209 Zeichen
\newline{}Handschrift: , deutsche Kurrent}\toendnotes[C]{\smallbreak}
\pstart
           \noindent{}{\pb}Lieber Guſtav, wenn Sie nichts luſtiges
      vorhaben und dazu aufgelegt ſind, machen
      Sie mir, der \label{K_L04186-1v}\edtext{Zi{\geminationm}erarreſt}{\lemma{\textnormal{\emph{Zimmerarrest}}}\Cendnote{\textnormal{Das vorliegende Korrespondenzstück ist
      undatiert. Es dürfte in die Zeit fallen, in denen Schnitzler wegen eines Gerstenkorns zuhause
         bleiben musste. Ein Besuch von Schwarzkopf\pwindex{Schwarzkopf, Gustav 7.\,11.\,1853 Wien – 13.\,11.\,1939 ebd.@\textsc{Schwarzkopf, Gustav} (7.\,11.\,1853 Wien – 13.\,11.\,1939 ebd.)|pwk} in diesen Tagen ist nicht im \emph{Tagebuch}\pwindex{Schnitzler, Arthur 15. 5. 1862 Wien – 21. 10. 1931 ebd.@\textsc{Schnitzler, Arthur} (15. 5. 1862 Wien – 21. 10. 1931 ebd.)!Tagebuch@\strich\emph{Tagebuch}|pwk}
         dokumentiert.}}}\label{K_L04186-1} hat, das
      Vergnügen, heut bei uns zu nachtmahlen.
      Je früher Sie kämen, je lieber.\pend
           \pstart Herzlichſt Ihr \spacefill\mbox{Arthur Schn.}\pend{}\selectlanguage{ngerman}\endnumbering\briefempfaengerindex{Schwarzkopf, Gustav@\textsc{Schwarzkopf, Gustav}!zzzSchnitzler, Arthur@\emph{von Arthur Schnitzler}!1897-12-112@{{[}zwischen 11. und 18. 12. 1897?{]}}|)be}\mylabel{L04186h}
\begin{anhang}
\end{anhang}\newcommand{\dateiname}{L04186}\newcommand{\titel}{Arthur Schnitzler an Gustav Schwarzkopf, [zwischen 11. und 18. 12. 1897?]}\newcommand{\editorInnen}{Herausgegeben von Jahnke, SelmaMüller, Martin Anton}\input{../tex-inputs/latex-leseansicht-abspann}
